\section{Decision two: which day takes the day}

\subsection{The table}

\rg{91} introduces it in a sentence that forecloses argument: ``Praecedentia
dierum liturgicorum, sublatis quibuslibet aliis titulis vel normis, unice
regitur ex sequenti \textsc{tabella dierum liturgicorum secundum ordinem
praecedentiae disposita}.'' Twenty-eight numbered rows, banded into the four
classes of day. The row numbers themselves are then used as addresses by a
dozen later rubrics --- ``diebus sub nn.~1, 2, 3 et 8 in tabella praecedentiae
recensitis'' --- so they must be learned as numbers, not merely as an order.

\begin{precedencetable}
1 & The feast of the Nativity of the Lord, Easter Sunday and Pentecost Sunday
(first class with an octave). & \\
2 & The Sacred Triduum. & \\
3 & The feasts of the Epiphany and the Ascension of the Lord, the Most Holy
Trinity, Corpus Christi, the Sacred Heart and Christ the King. & \\
4 & The feasts of the Immaculate Conception and the Assumption of the Blessed
Virgin Mary. & \\
5 & The vigil and the octave day of the Nativity. & \\
6 & The Sundays of Advent, Lent and Passiontide, and Low Sunday. & \\
7 & The first-class ferias not named above: Ash Wednesday and Monday, Tuesday
and Wednesday of Holy Week. & \\
8 & The Commemoration of All the Faithful Departed --- \emph{which
nevertheless yields to an occurring Sunday}. & The only row of the table
carrying its own defeat.\\
9 & The vigil of Pentecost. & \\
10 & The days within the octaves of Easter and Pentecost. & \\
11 & First-class feasts of the universal Church not named above. & Addressed
by name in \rgm{406}\,b.\\
12 & Proper first-class feasts: principal Patron of nation, of region or
province, of diocese; anniversary of the dedication of the cathedral;
principal Patron of the place, town or city; feast and anniversary of the
dedication of one's own church, or of a public or semi-public oratory taking a
church's place; Title of one's own church; Title of an Order or Congregation;
canonised Founder; principal Patron of an Order or Congregation and of a
religious province. & Order within the row as printed.\\
13 & Indulted first-class feasts, movable first, then fixed. & \\
14 & Second-class feasts of the Lord, movable first, then fixed. & \\
15 & Second-class Sundays. & \\
16 & Second-class feasts of the universal Church which are not of the Lord. & \\
17 & The days within the octave of the Nativity. & \\
18 & Second-class ferias: Advent from 17 to 23 December inclusive, and the
Ember Days of Advent, Lent and September. & \\
19 & Proper second-class feasts: secondary Patron of nation, region or
province, diocese, place, town or city; feasts of saints or blessed under
\rg{43}\,d; feasts of saints proper to a church (\rg{45}\,c); beatified
Founder (\rg{46}\,b); secondary Patron of an Order or Congregation and of a
religious province (\rg{46}\,d); feasts of saints or blessed under
\rg{46}\,e. & \\
20 & Indulted second-class feasts, movable first, then fixed. & \\
21 & Second-class vigils. & \\
22 & The ferias of Lent and Passiontide from the Thursday after Ash Wednesday
to the Saturday before Passion Sunday inclusive, Ember Days excepted. &
\textbf{Above} rows 23 and 24.\\
23 & Third-class feasts inscribed in particular calendars: first the proper
ones (saints or blessed under \rg{43}\,d; feasts of blessed proper to a church,
\rg{45}\,d; saints or blessed under \rg{46}\,e), then indulted ones, movable
first, then fixed. & \\
24 & Third-class feasts inscribed in the calendar of the universal Church,
movable first, then fixed. & \\
25 & The ferias of Advent to 16 December inclusive, Ember Days excepted. &
\textbf{Below} rows 23 and 24.\\
26 & Third-class vigils. & \\
27 & The Office of Our Lady on Saturday. & \\
28 & Fourth-class ferias. & \\
\end{precedencetable}

Three features of the table are load-bearing and easy to miss.

\begin{itemize}
\item \textbf{It ranks days, not saints.} Rows 6, 7, 15, 18, 22, 25, 27 and 28
contain no feast at all. A saint whose feast is third class in the universal
calendar sits at row 24, below the Lenten ferias at row 22 and below any
third-class feast a particular calendar owns at row 23 --- including a
particular calendar's indulted feasts.

\item \textbf{Universal beats particular within a class, except at third
class.} Rows 11 and 12 put universal first-class feasts above proper ones; rows
16 and 19 do the same at second class. At third class the order reverses: row
23 (particular) precedes row 24 (universal). The code does not explain the
inversion; it simply prints it, and the practical consequence is that a
diocesan third-class feast displaces a universal one on the same date.

\item \textbf{Row 8 contains a defeat clause.} All Souls is a first-class
liturgical day placed above the vigil of Pentecost and above every first-class
feast of the universal Church, yet the row itself adds ``quae tamen locum
cedit dominicae occurrenti.'' \rg{16}\,b says the same thing from the Sunday's
side. This is the only place in the table where a row is qualified in this
way, and it drives Case~3 below.
\end{itemize}

\subsection{Occurrence: accidental and perpetual}

\begin{rulebox}{\rg{92}}
Occurrentia dicitur occursus duorum vel plurium Officiorum uno eodemque die.
Occurrentia autem dicitur \textit{accidentalis}, quando dies liturgicus
mobilis et dies liturgicus fixus certis solummodo annorum intervallis simul
occurrunt; \textit{perpetua} vero quando duo dies liturgici quotannis simul
occurrunt.
\end{rulebox}

The distinction is not descriptive. It selects which remedy is available.
Accidental occurrence --- a movable day meeting a fixed one only in some years
--- may lead to \term{transfer} (\rg{95}--\rg{99}). Perpetual occurrence --- two
days that meet every year --- may lead to \term{reposition}, a permanent
reassignment in the calendar (\rg{100}--\rg{102}). A feast cannot be
``transferred'' out of a perpetual collision, and cannot be ``repositioned''
out of an accidental one.

\rg{93} states the general effect and enumerates the only four outcomes:

\begin{rulebox}{\rg{93}}
Effectus occurrentiae est, ut Officium diei liturgici gradus inferioris
Officio gradus superioris cedat: quod fieri potest per minus nobilis
omissionem, aut commemorationem, aut translationem, aut repositionem prout
numeris sequentibus indicatur.
\end{rulebox}

\noindent
Omission, commemoration, transfer, reposition. Section~4 takes them in turn.
Note in passing that \rg{93} speaks of the \emph{Office}: precedence decides
what the Office of the day is, and the Mass follows the Office only because
\rgm{270} separately says so (``Missa proinde per se cum Officio diei
concordare debet''), with its own list of exceptions.

One special case is settled here rather than later. \rg{95} closes with a rule
about identity rather than rank:

\begin{rulebox}{\rg{95} (second paragraph)}
Si vero duo festa eiusdem Divinae Personae aut duo festa eiusdem Sancti vel
Beati simul occurrunt, fit de festo, quod in tabella praecedentiae superiorem
obtinet locum, et aliud omittitur.
\end{rulebox}

\noindent
Two feasts of the same divine Person, or of the same saint or blessed, do not
commemorate one another and the loser is not transferred: it is simply
omitted. \rg{112}\,a extends the same principle to commemorations, and
\rgm{317} extends it to votive Masses.

\subsection{Concurrence, and what it does not decide}

\begin{rulebox}{\rg{103}--\rg{105}}
Concurrentia dicitur concursus Vesperarum diei liturgici currentis cum I
Vesperis diei liturgici subsequentis. \dots{} In concurrentia praeferuntur
Vesperae diei liturgici classis superioris, et alterae commemorantur vel non,
iuxta rubricas. Quando vero dies liturgici, quorum Vesperae concurrunt, sunt
eiusdem classis, dicuntur integrae secundae Vesperae de Officio currenti et
fit commemoratio sequentis, iuxta rubricas.
\end{rulebox}

Concurrence disposes of one evening's Vespers. It never decides which Mass is
said the next morning; that is occurrence, settled by the same table but on
the following day's own competitors. The confusion is easy to fall into
because \rg{106} says the commemoration rules ``valent tam pro Missa quam pro
Officio, sive in occurrentia sive in concurrentia'' --- the same
\emph{commemoration} apparatus serves both, which is not the same as the two
questions being one question.

\begin{caution}{What the 1962 Missal does not print}
The promulgated code in the \work{Acta} carries, after the enumerations of
days, a \textsc{tabella occurrentiae}, a \textsc{tabella concurrentiae} and a
short set of \textsc{notanda} explaining both. The 1962 Vatican typical
edition of the Missal reprints the precedence table of \rg{91} and does not
reprint the occurrence or concurrence tables or the notanda. A reader working
from the Missal alone therefore has the ranking table but not the
Acta's tabular shortcut for the outcome, and must reason from
\rg{92}--\rg{114}. Nothing is lost in substance --- the notanda restate
\rg{16}\,a, \rg{95} and a Vespers rule --- but the omission is worth knowing
before one goes looking for a table that is not there.
\end{caution}
